% capitulos/1planteamiento.tex
\section{Planteamiento del Problema}

\subsection{Identificaci\'on del Problema}

Los sistemas de detecci\'on autom\'atica de personas son componentes
cr\'{i}ticos en aplicaciones de seguridad p\'{u}blica, monitoreo
industrial, rob\'otica m\'ovil y asistencia m\'edica, entre otras.
En estos contextos, los detectores modernos de una y dos etapas
---como YOLO~\cite{Redmon2016} y Faster R-CNN~\cite{Ren2015}---
logran resultados destacados bajo condiciones controladas de
iluminaci\'on y visibilidad. Sin embargo, su efectividad disminuye
dr\'asticamente al enfrentarse a dos factores adversos: la
\textbf{oclusi\'on} ---obstrucci\'on parcial o total de la silueta
corporal por otro objeto, persona o elemento del entorno--- y la
\textbf{baja visibilidad} ---causada por escasa iluminaci\'on
nocturna, niebla, humo o lluvia.

La brecha radica en que los enfoques actuales abordan estas dos
limitaciones visuales de forma \textbf{fragmentada e independiente}.
Los modelos optimizados para baja iluminaci\'on no incorporan
mecanismos robustos de manejo de oclusi\'on, y viceversa. Weng \&
Niu~\cite{ElsYolo2025} demostraron que arquitecturas basadas en
YOLOv11s pueden alcanzar un mAP@0.5 de 74{,}3\,\% en condiciones
de baja visibilidad sobre el dataset ExDark, pero su modelo no
integra ning\'un m\'odulo espec\'{i}fico para oclusi\'on parcial o
severa. Esta carencia se traduce en ca\'{i}das significativas de
rendimiento cuando la silueta humana queda parcialmente cubierta,
situaci\'on com\'un en entornos reales como espacios concurridos,
pasillos o escenas nocturnas con obst\'aculos.

Actualmente no existe una arquitectura unificada basada en
YOLOv11 que gestione simult\'aneamente la baja visibilidad y la
oclusi\'on severa en im\'agenes RGB convencionales, sin requerir
sensores especializados como c\'amaras t\'ermicas o de infrarrojo.
Esto limita la aplicabilidad de los detectores existentes en
entornos reales donde se emplean c\'amaras est\'andar~\cite{Deng2024,
Tang2026, He2025}.

\bigskip
\noindent\textbf{Pregunta de investigaci\'on:}

\begin{quote}
\textit{?`En qu\'e medida el dise\~{n}o del modelo ARGOS-11, basado
en YOLOv11s con un m\'odulo de atenci\'on cruzada y augmentaci\'on
sint\'etica de oclusi\'on, mejora la precisi\'on media (mAP@0.5)
en la detecci\'on de personas bajo condiciones simult\'aneas de
oclusi\'on severa y baja visibilidad en im\'agenes RGB, respecto al
modelo de referencia ELS-YOLO?}
\end{quote}

%---------------------------------------------------------------
\subsection{Justificaci\'on}
%---------------------------------------------------------------

\textbf{Justificaci\'on te\'orica.}
La literatura cient\'{i}fica reciente confirma que los modelos de
detecci\'on de personas publicados entre 2022 y 2026 abordan la
oclusi\'on y la baja visibilidad como problemas independientes.
Ning\'un trabajo propone una arquitectura unificada que gestione
ambas condiciones simult\'aneamente en im\'agenes RGB
convencionales. El modelo ARGOS-11 aportar\'a un nuevo enfoque
te\'orico que integra mecanismos de atenci\'on cruzada para
oclusi\'on sobre la base de un detector YOLOv11s optimizado para
baja iluminaci\'on, contribuyendo al avance del estado del arte en
detecci\'on robusta de personas~\cite{ElsYolo2025, Deng2024}.

\textbf{Justificaci\'on pr\'actica.}
La detecci\'on de personas en condiciones adversas es un requisito
exigente en m\'ultiples dominios: vigilancia de seguridad p\'ublica,
monitoreo de obras de construcci\'on, sistemas de asistencia para
personas mayores, rob\'otica de rescate y an\'alisis de flujo
peatonal nocturno. Los detectores actuales reportan ca\'{i}das
severas de rendimiento cuando la silueta humana se obstruye m\'as
del 70\,\% en escenas de baja iluminaci\'on~\cite{Kim2025, Tang2026}.
ARGOS-11 busca resolver esta limitaci\'on trabajando exclusivamente
con im\'agenes RGB est\'andar, sin requerir hardware especializado,
lo que lo hace deployable en cualquier c\'amara convencional.

\textbf{Justificaci\'on metodol\'ogica.}
Esta investigaci\'on aplica el pipeline de entrenamiento
establecido por Weng \& Niu~\cite{ElsYolo2025} en ELS-YOLO
---optimizer SGD, 200 \'epocas, batch size 16, dataset ExDark---
como base reproducible y comparable. A este pipeline se agregan dos
contribuciones metodol\'ogicas propias: (1) un m\'odulo de
atenci\'on cruzada canal-espacial para manejo de oclusi\'on, y
(2) una estrategia de augmentaci\'on sint\'etica de oclusi\'on
durante el entrenamiento. Esta combinaci\'on permite una
comparaci\'on directa y justa con el modelo de referencia, dado
que ambos comparten el mismo dataset, optimizador y
configuraci\'on de entrenamiento.

%---------------------------------------------------------------
\subsection{Estado del Arte}
%---------------------------------------------------------------

La detecci\'on autom\'atica de personas ha evolucionado desde los
detectores cl\'asicos basados en histogramas de gradientes
orientados (HOG) hasta las arquitecturas profundas modernas.
Doll\'ar et al.~\cite{Dollar2012} demostraron que la oclusi\'on es
el factor que mayor degradaci\'on produce en el rendimiento de los
detectores. Con el aprendizaje profundo, Ren et al.~\cite{Ren2015}
introdujeron Faster R-CNN, mientras que Redmon et
al.~\cite{Redmon2016} propusieron YOLO como el paradigma de
detecci\'on en tiempo real.

Los trabajos m\'as relevantes al presente proyecto son:

\begin{enumerate}

\item \textbf{ELS-YOLO~\cite{ElsYolo2025} (referencia metodol\'ogica
principal).} Weng \& Niu propusieron ELS-YOLO, un modelo ligero
basado en YOLOv11s dise\~{n}ado para condiciones de baja
visibilidad. El modelo incorpora un backbone re-parametrizado
(ER-HGNetV2) con mecanismos de atenci\'on de canal eficiente
(ECA), una pir\'amide de selecci\'on de caracter\'{i}sticas
ligera (LFSPN) y un cabezal con normalizaci\'on por lotes
compartida (SCSHead). Evaluado en el dataset ExDark, alcanza un
\textbf{mAP@0.5 de 74{,}3\,\%} manteniendo inferencia en tiempo
real. Este resultado constituye la l\'{i}nea de base que
ARGOS-11 debe superar.

\item \textbf{PYOLO~\cite{Hu2025}.} Hu et al. propusieron PYOLO,
un modelo espec\'{i}ficamente dise\~{n}ado para oclusi\'on
peatonal. Incorpora m\'odulos de predicci\'on de partes
corporales visibles que permiten inferir la posici\'on global de
la persona a partir de regiones parcialmente expuestas (cabeza,
abdomen). Sus resultados sobre Caltech demuestran mejoras
consistentes ante oclusi\'on moderada y severa.

\item \textbf{EAFF-Net~\cite{Shen2025}.} Shen et al. presentaron
una red de fusi\'on dual con atenci\'on eficiente para detecci\'on
peatonal bimodal. El m\'odulo EAFF (Efficient Attention Feature
Fusion) sirve de inspiraci\'on para el dise\~{n}o del m\'odulo de
atenci\'on cruzada de ARGOS-11, adaptado para im\'agenes RGB
\'unicas en lugar de pares multimodales.

\item \textbf{Deng et al.~\cite{Deng2024}.} Propusieron un detector
multiespectral ligero con atenci\'on espacial re-ponderada. Sus
experimentos en KAIST confirman que la atenci\'on cruzada
canal-espacial mejora la detecci\'on en escenas con baja
iluminaci\'on, reportando mejoras de hasta 8 puntos en mAP
respecto a detectores de referencia.

\item \textbf{Tang et al.~\cite{Tang2026}.} Propusieron un
razonamiento jer\'arquico adaptativo para oclusi\'on en entornos
din\'amicos con recursos computacionales limitados. Sus
experimentos confirman que los detectores actuales presentan
ca\'{i}das significativas cuando la oclusi\'on supera el 70\,\%
de la silueta, brecha que ARGOS-11 busca mitigar.

\end{enumerate}

El an\'alisis del estado del arte revela que ning\'un trabajo
combina un detector YOLOv11 con atenci\'on cruzada para oclusi\'on
evaluado sobre im\'agenes RGB de baja visibilidad. Esta es la
brecha espec\'{i}fica que ARGOS-11 aborda.

%---------------------------------------------------------------
\subsection{Objetivos}
%---------------------------------------------------------------

\subsubsection{Objetivo General}

Desarrollar el modelo ARGOS-11, basado en YOLOv11s con un m\'odulo
de atenci\'on cruzada canal-espacial y augmentaci\'on sint\'etica
de oclusi\'on, que alcance una precisi\'on media (mAP@0.5) igual
o superior al 74{,}3\,\% del modelo de referencia ELS-YOLO en la
detecci\'on de personas bajo condiciones simult\'aneas de oclusi\'on
severa y baja visibilidad en im\'agenes RGB convencionales.

\subsubsection{Objetivos Espec\'{i}ficos}

\begin{enumerate}[label=OE\arabic*.]

\item Dise\~{n}ar e integrar un m\'odulo de atenci\'on cruzada
      canal-espacial en la arquitectura YOLOv11s, inspirado en el
      mecanismo EAFF de Shen et al.~\cite{Shen2025} y adaptado
      para im\'agenes RGB en condiciones de baja iluminaci\'on y
      oclusi\'on simult\'aneas.

\item Implementar una estrategia de augmentaci\'on sint\'etica de
      oclusi\'on sobre el dataset ExDark, cubriendo entre el 30\,\%
      y el 70\,\% de las siluetas de personas con parches opacos
      aleatorios durante el entrenamiento, para incrementar la
      robustez del modelo ante oclusi\'on severa.

\item Entrenar el modelo ARGOS-11 siguiendo el pipeline de Weng \&
      Niu~\cite{ElsYolo2025}: optimizador SGD con tasa de
      aprendizaje inicial de 0{,}01, momentum de 0{,}937, weight
      decay de 0{,}0005, 200 \'epocas y batch size de 16, sin pesos
      preentrenados, sobre el dataset ExDark.

\item Evaluar el modelo ARGOS-11 midiendo mAP@0.5, mAP@0.5:0.95 y
      FPS sobre el conjunto de prueba de ExDark, comparando los
      resultados con ELS-YOLO como baseline, para demostrar que
      ARGOS-11 es igual o superior en precisi\'on manteniendo
      capacidad de inferencia en tiempo real.

\end{enumerate}

%---------------------------------------------------------------
\subsection{Hip\'otesis}
%---------------------------------------------------------------

\textbf{Hip\'otesis de trabajo (H1):}

El modelo ARGOS-11, que integra un m\'odulo de atenci\'on cruzada
canal-espacial y augmentaci\'on sint\'etica de oclusi\'on sobre la
arquitectura YOLOv11s, alcanza una precisi\'on media (mAP@0.5)
igual o superior al 74{,}3\,\% obtenido por ELS-YOLO en el dataset
ExDark, bajo condiciones simult\'aneas de oclusi\'on severa (entre
el 30\,\% y el 70\,\% de la silueta obstruida) y baja visibilidad
(escenas nocturnas, interiores oscuros, niebla y lluvia), trabajando
exclusivamente con im\'agenes RGB convencionales.

\bigskip
\textbf{Hip\'otesis nula (H0):}

El modelo ARGOS-11 obtiene una precisi\'on media (mAP@0.5) inferior
al 74{,}3\,\% de ELS-YOLO en el dataset ExDark bajo las mismas
condiciones de oclusi\'on severa y baja visibilidad, indicando que
las modificaciones propuestas no aportan mejora significativa
respecto al modelo de referencia.

\bigskip
\textbf{Sustento bibliogr\'afico de H1:}

\begin{itemize}
\item \textbf{Weng \& Niu~(2025)~\cite{ElsYolo2025}} establecen
      que YOLOv11s, con modificaciones arquitect\'onicas ligeras,
      puede alcanzar mAP@0.5 de 74{,}3\,\% en baja visibilidad sobre
      ExDark manteniendo tiempo real. Nuestro m\'odulo de atenci\'on
      adicional deber\'{i}a mejorar ese resultado al agregar
      robustez ante oclusi\'on, sin sacrificar velocidad de
      inferencia.

\item \textbf{Shen et al.~(2025)~\cite{Shen2025}} demuestran que
      los mecanismos de atenci\'on cruzada canal-espacial mejoran
      significativamente la localizaci\'on de regiones parcialmente
      ocluidas. Este principio, aplicado a RGB, justifica la mejora
      esperada en mAP de ARGOS-11 sobre ELS-YOLO.

\item \textbf{Deng et al.~(2024)~\cite{Deng2024}} reportan mejoras
      de hasta 8 puntos en mAP al agregar atenci\'on espacial
      cruzada sobre detectores base en escenas de baja iluminaci\'on,
      respaldando directamente la hip\'otesis de que ARGOS-11
      superar\'a el 74{,}3\,\% de la l\'{i}nea base.
\end{itemize}

%---------------------------------------------------------------
\subsection{Variables}
%---------------------------------------------------------------

A continuaci\'on se presenta la operacionalizaci\'on de las
variables que intervienen en el estudio.

\bigskip

\begin{table}[h!]
\centering
\caption{Tabla de operacionalizaci\'on de variables de ARGOS-11}
\label{tab:variables}
\renewcommand{\arraystretch}{1.4}
\begin{tabular}{|p{3cm}|c|p{3.3cm}|p{3.3cm}|p{1.6cm}|}
\hline
\textbf{Variable} & \textbf{Rol} &
\textbf{Definici\'on conceptual} &
\textbf{Definici\'on operacional} & \textbf{Escala} \\
\hline
Arquitectura del modelo (ARGOS-11 vs.\ ELS-YOLO) &
VI &
Configuraci\'on del detector: YOLOv11s base (ELS-YOLO) versus
YOLOv11s con m\'odulo de atenci\'on cruzada y augmentaci\'on de
oclusi\'on (ARGOS-11). &
Presencia o ausencia del m\'odulo de atenci\'on cruzada y la
augmentaci\'on sint\'etica de oclusi\'on en el pipeline de
entrenamiento. &
Nominal (Discreta) \\
\hline
Precisi\'on Media (mAP@0.5) &
VD1 &
M\'etrica est\'andar de exactitud del detector en detecci\'on de
objetos al umbral IoU = 0{,}5. &
C\'alculo del \'area bajo la curva Precision-Recall evaluado sobre
el conjunto de prueba del dataset ExDark, umbral IoU = 0{,}5. &
Raz\'on (Continua) \\
\hline
Fotogramas por Segundo (FPS) &
VD2 &
Tasa de procesamiento del modelo en tiempo real. &
N\'umero de frames procesados correctamente por segundo durante
inferencia sobre hardware de referencia (GPU Jetson Nano /
RTX 5070 Ti). &
Raz\'on (Continua) \\
\hline
Nivel de oclusi\'on sint\'etica &
VC &
Grado de obstrucci\'on aplicado sobre las anotaciones de personas
en ExDark durante el entrenamiento. &
Porcentaje del bounding box de cada persona cubierto por parches
opacos aleatorios durante la augmentaci\'on de entrenamiento. &
Ordinal (Discreta) \\
\hline
\end{tabular}
\end{table}

\bigskip
\newpage

\noindent\textbf{Variable independiente (VI):} Arquitectura del
modelo --- YOLOv11s con m\'odulo de atenci\'on cruzada
canal-espacial y augmentaci\'on sint\'etica de oclusi\'on
(ARGOS-11) versus YOLOv11s base sin modificaciones (ELS-YOLO).
Instrumento: c\'odigo fuente del modelo (Python/PyTorch).
Valores: \{ARGOS-11, ELS-YOLO baseline\}.

\vspace{0.5cm}
\noindent\textbf{Variable dependiente 1 (VD1):} Precisi\'on Media
(mAP@0.5) --- \'area bajo la curva Precision-Recall al umbral
IoU = 0{,}5. Instrumento: COCO evaluation API (pycocotools).
Meta: $\geq 74{,}3\,\%$ (resultado de referencia de ELS-YOLO).

\vspace{0.5cm}
\noindent\textbf{Variable dependiente 2 (VD2):} Fotogramas por
Segundo (FPS) --- cantidad de frames procesados por segundo.
Instrumento: profiler de GPU. Umbral m\'{i}nimo: 30 FPS para
viabilidad en tiempo real.

\vspace{0.5cm}
\noindent\textbf{Variable de control (VC):} Nivel de oclusi\'on
sint\'etica --- porcentaje de la silueta cubierta durante
augmentaci\'on. Valores: \{Ligera ($\leq$30\,\%), Moderada
(31--69\,\%), Severa ($\geq$70\,\%)\}. Instrumento: script de
augmentaci\'on propio implementado en Python/Albumentations.
